% Use only LaTeX2e, calling the article.cls class and 12-point type.

\documentclass[12pt]{article}

% Users of the {thebibliography} environment or BibTeX should use the
% scicite.sty package, downloadable from *Science* at
% www.sciencemag.org/about/authors/prep/TeX_help/ .
% This package should properly format in-text
% reference calls and reference-list numbers.

\usepackage{scicite}

% Use times if you have the font installed; otherwise, comment out the
% following line.

\usepackage{url}
\usepackage{times}

% The preamble here sets up a lot of new/revised commands and
% environments.  It's annoying, but please do *not* try to strip these
% out into a separate .sty file (which could lead to the loss of some
% information when we convert the file to other formats).  Instead, keep
% them in the preamble of your main LaTeX source file.


%For the long table of political userhandles:
\usepackage{booktabs}
\usepackage{longtable}
\usepackage{array}


% The following parameters seem to provide a reasonable page setup.

\topmargin 0.0cm
\oddsidemargin 0.2cm
\textwidth 16cm 
\textheight 21cm
\footskip 1.0cm

% Since this is an initial submission, it is permitted to use the command "\includegraphics{}". Therefore, the following package is needed:
\usepackage{graphicx}

%The next command sets up an environment for the abstract to your paper.

\newenvironment{sciabstract}{%
\begin{quote} \bf}
{\end{quote}}


% If your reference list includes text notes as well as references,
% include the following line; otherwise, comment it out.

\renewcommand\refname{References and Notes}

% The following lines set up an environment for the last note in the
% reference list, which commonly includes acknowledgments of funding,
% help, etc.  It's intended for users of BibTeX or the {thebibliography}
% environment.  Users who are hand-coding their references at the end
% using a list environment such as {enumerate} can simply add another
% item at the end, and it will be numbered automatically.

\newcounter{lastnote}
\newenvironment{scilastnote}{%
\setcounter{lastnote}{\value{enumiv}}%
\addtocounter{lastnote}{+1}%
\begin{list}%
{\arabic{lastnote}.}
{\setlength{\leftmargin}{.22in}}
{\setlength{\labelsep}{.5em}}}
{\end{list}}


% Include your paper's title here

\title{Digital literacy reduces susceptibility to populist content on social media} 

% Place the author information here.  Please hand-code the contact
% information and notecalls; do *not* use \footnote commands.  Let the
% author contact information appear immediately below the author names
% as shown.  We would also prefer that you don't change the type-size
% settings shown here.

\author
{Kristina Aleksandrovna Pedersen$^{1}$, Ferdinand Geissler$^{2}$, Felix Hartmann$^{2}$, \\Johannes Giesecke$^{1}$, Macartan Humphreys$^{3}$, Heike Klüver$^{2\ast}$\\
\\
\normalsize{$^{1}$ Department of International Economics, Government and Business,}\\ 
\normalsize{$^{2}$Department of Social Sciences, Humboldt University of Berlin,}\\
\normalsize{Unter den Linden 6, 10099 Berlin, Germany}\\
\normalsize{Copenhagen Business School, Solbjerg Plads 3, 2000 Frederiksberg, Denmark}\\
\normalsize{$^{3}$WZB Berlin Social Science Center,}\\
\normalsize{Reichpietschufer 50, 10785 Berlin, Germany}\\
\\
\normalsize{$^\ast$To whom correspondence should be addressed; E-mail:  
heike.kluever@hu-berlin.de}
}

% Include the date command, but leave its argument blank.

\date{}



%%%%%%%%%%%%%%%%% END OF PREAMBLE %%%%%%%%%%%%%%%%



\begin{document} 

% Double-space the manuscript.

\baselineskip24pt

% Make the title.

\maketitle 



% Place your abstract within the special {sciabstract} environment.

%\begin{sciabstract}
%  Populist actors increasingly use social media to disseminate emotionally manipulative political messages. Can digital literacy training reduce citizens’ susceptibility to such content? We address this question with a large-scale online experiment conducted during the 2025 German federal election. Over 4,400 participants were randomly assigned to a control group or to a short, interactive digital literacy workshop designed to teach strategies for recognizing emotional manipulation in political TikToks. Participants then viewed a randomized sequence of authentic TikTok videos from German political parties and entertainment accounts. We find that workshop participants were significantly less likely to like and more likely to report emotionally manipulative political posts, particularly those disseminated by populist parties. These effects are strongest for content invoking anger or superiority and for posts from parties participants supported. Moreover, the workshop increased watch time and reduced skipping, suggesting greater cognitive engagement. Our findings demonstrate that brief digital interventions can mitigate the persuasive power of emotionally manipulative populist content on social media.
%\end{sciabstract}


\begin{sciabstract}
Populists increasingly use short-form video platforms to disseminate emotionally manipulative political messages. Can digital literacy interventions reduce citizens’ engagement with such content? We conducted an online experiment during the 2025 German election. Over 4,400 participants were randomly assigned to a control or a short, interactive workshop that taught how to recognize emotional manipulation. Participants then interacted with a TikTok-like feed containing authentic videos from party and entertainment accounts. The workshop reduced liking and increased reporting of emotionally manipulative political posts, with the largest effects for posts from populist parties. Treated participants also watched longer and skipped less, consistent with more deliberate processing. These results show that a brief, scalable intervention can alter behavioral responses to emotionally manipulative political content in a realistic short-video environment.
\end{sciabstract}


% Word count: 2884 (without abstract)


% In setting up this template for *Science* papers, we've used both
% the \section* command and the \paragraph* command for topical
% divisions.  Which you use will of course depend on the type of paper
% you're writing.  Review Articles tend to have displayed headings, for
% which \section* is more appropriate; Research Articles, when they have
% formal topical divisions at all, tend to signal them with bold text
% that runs into the paragraph, for which \paragraph* is the right
% choice.  Either way, use the asterisk (*) modifier, as shown, to
% suppress numbering.

% =========================
% REPLACEMENT: INTRODUCTION (tightened + sharper contribution)
% =========================
\section*{Introduction}

As political communication migrates to short-form, algorithmically curated video platforms, populist actors increasingly rely on emotionally charged and manipulative cues to capture attention and shape engagement. Whether citizens can be equipped to recognize and resist these cues at the moment of exposure is an open question.

Evidence on media and digital literacy interventions is mixed. Several studies report improved discernment, whereas others find limited or null effects (e.g., \citen{AmarEtAl2025_MediaLiteracyIndia,BADRINATHAN_2021,carletonathey2023,kozyreva2024,soetekouw2024, guess2020digital,kim_going_2024, gonzalez-aguilar_populist_2023}). Beyond heterogeneous findings, most existing work (i) focuses on static, text-based misinformation; (ii) relies on survey or hypothetical judgments that assess whether people \textit{can} detect manipulation when prompted rather than whether they \textit{change behavior} when exposed; and (iii) rarely examines emotionally manipulative political content embedded in fast-paced video feeds. 

This gap is consequential because algorithmic ranking systems on short-video platforms tend to reward content that maximizes engagement, and emotionally charged political posts reliably attract attention. Platform and electoral incentives thus converge: the same emotional appeals that boost engagement also advance political actors’ goals of visibility and mobilization, producing a reinforcing dynamic that can privilege manipulative content over deliberative discourse (\citen{milli_engagement_2025, piccardi_reranking_2025}). Yet we know little about whether brief training interventions can change how citizens engage with emotionally manipulative political content in the very environments where such messages now increasingly shape public discourse.

We address this gap with a preregistered online experiment conducted with 4,437 participants during the 2025 German federal election. Participants were randomly assigned to a control group or to a brief digital literacy workshop designed to help them recognize emotional manipulation (fear, anger, superiority) and apply simple “stop-and-think” strategies. The intervention draws on inoculation theory: exposure to weakened examples of manipulative techniques paired with explanations can foster “cognitive antibodies” that increase resistance to subsequent persuasive attempts \cite{mcguire1961resistance,amazeen_science}.

After the intervention, participants interacted for five minutes with a custom-built platform that replicated core TikTok affordances (sequential short-video exposure and standard engagement buttons). The feed contained real TikTok videos from German party accounts posted during the election campaign as well as non-political entertainment content. We unobtrusively measured behavioral outcomes (liking, sharing, reporting, bookmarking, skipping, and watch time) to estimate causal effects on engagement in a realistic short-video environment.

We find that the workshop reduced liking and increased reporting of emotionally manipulative political videos, particularly those invoking anger or superiority and those disseminated by populist parties. Effects were largest when videos came from a respondent’s preferred party. Treated participants also skipped less and watched longer, suggesting more deliberate processing rather than simple avoidance. Together, these results provide behavioral evidence that a brief, scalable intervention can change how people respond to emotionally manipulative political content in short-form video feeds.



%\section*{Introduction}

%As political communication migrates to short-form, algorithmically curated video platforms, populist actors increasingly rely on emotionally manipulative cues to mobilize attention and engagement. Can brief digital-literacy training help people recognize and resist such content when they encounter it?

%Evidence on literacy interventions is mixed: several studies report improved discernment, but others find limited or null effects, including a prominent field campaign in India (e.g., \citen{AmarEtAl2025_MediaLiteracyIndia,BADRINATHAN_2021,carletonathey2023,kozyreva2024,soetekouw2024, guess2020digital,kim_going_2024, gonzalez-aguilar_populist_2023}). Beyond mixed results, most existing work (i) focuses on static, text-based misinformation; (ii) relies on hypothetical or survey outcomes that test whether people \textit{can} spot manipulation when prompted rather than whether they \textit{do} behave differently when exposed; and (iii) rarely examines emotionally manipulative \textit{political} content in fast-paced video feeds. In short, the literature diverges on effectiveness and has limited evidence from realistic, video-based contexts where manipulation is increasingly delivered.

%We address this gap with a large-scale online experiment conducted among more then 4,400 participants during the 2025 German federal election. We test whether a brief, digital literacy workshop online can reduce engagement with emotionally manipulative political TikTok content. The workshop teaches participants to identify emotional triggers - such as fear, anger, or superiority - and apply simple “stop-and-think” strategies when confronted with emotionally provocative posts. Our intervention is based on inoculation theory. The central idea is that, much like exposure to a virus triggers the production of antibodies, exposure to weakened forms of manipulative content together with explanations of the underlying techniques can foster “cognitive antibodies” that increase individuals’ resistance to politically manipulative content on social media \cite{mcguire1961resistance,amazeen_science}. 

%After the intervention, participants were shown a randomized feed of real TikTok videos, including content from entertainment accounts and official political party accounts that had posted during the election campaign.  We measured their engagement behavior - liking, sharing, reporting, bookmarking, skipping, and watch time - allowing us to assess the causal impact of the intervention in a realistic digital environment.

%We find that the workshop significantly reduced likes and increased reporting of emotionally manipulative political videos, particularly those invoking anger or superiority and those disseminated by populist parties. These effects were strongest when the content aligned with participants’ political preferences. Participants in the treatment group also spent more time watching and were less likely to skip both political and entertainment videos, suggesting greater cognitive engagement.

%These results provide new evidence that digital literacy can work in a high-stakes, polarized setting, and in the very medium where manipulative appeals are delivered: short-form political video. As political communication continues to migrate to algorithmic platforms, fostering emotional awareness and digital resilience may be critical to safeguarding democratic discourse.


%%%%%%%%%%%%% REVISED STRUcTURE

\section*{Results}

\subsection*{Study design}
We conducted a preregistered online experiment during the 2025 German federal election ($N=4{,}437$) to test whether a brief digital literacy workshop alters behavioral engagement with emotionally manipulative political content in short-form video feeds. Respondents were randomly assigned to a control condition or to a short, interactive workshop designed to help them recognize emotional manipulation (fear, anger, superiority) and apply a simple stop-and-think routine (for more details, see Materials and Methods). Immediately afterward, participants interacted for five minutes with a custom TikTok-like feed containing authentic political videos from German party accounts as well as non-political entertainment content. We unobtrusively recorded behavioral engagement (liking, sharing, reporting, bookmarking) as well as attention measures (skipping and watch time). 
Open-ended responses collected after the browsing task suggest that respondents largely treated the platform as real rather than as part of an experiment (SM, section~\ref{appendix:ecological_validity}).

\subsection*{Primary outcomes: engagement with political TikTok videos}

We begin by examining treatment effects on engagement with emotionally manipulative political videos. Figure~\ref{fig:overall-patterns} shows that the workshop significantly reduced \textit{liking} of emotionally manipulative political posts by approximately one percentage point relative to the control group. Given a baseline liking rate of about nine percent, this corresponds to a meaningful relative reduction. Liking of neutral political and entertainment content remained largely unchanged.

In contrast, we observe a clear increase in \textit{reporting} behavior. Treated respondents were approximately 1.3 percentage points more likely to report emotionally manipulative political videos than respondents in the control condition. Given a reporting baseline of roughly two percent, this represents a substantial increase. Reporting of neutral political or entertainment content did not increase in a comparable magnitude.

By comparison, treatment effects on \textit{sharing} are small and not statistically distinguishable from zero. Although point estimates suggest a slight decline in sharing of emotionally manipulative political videos, confidence intervals include zero.

We observe a modest increase in \textit{bookmarking} of political videos (about 0.6 percentage points). Because bookmarking may reflect heterogeneous motivations (e.g., saving content for later viewing rather than endorsement), we interpret this outcome cautiously. Overall, the intervention primarily reduced positive reinforcement (likes) and increased governance-oriented responses (reports) toward emotionally manipulative political content.

\subsection*{Heterogeneity by party and partisan alignment}

Treatment effects vary across political senders (Figure~\ref{fig:patterns-party}). The largest reduction in liking is observed for emotionally manipulative posts disseminated by the right-wing populist \emph{AfD}. Treated respondents were more than two percentage points less likely to like AfD’s manipulative content relative to controls. We also observe a significant reduction for the left-wing populist \textit{Die Linke}, whereas effects for other parties are smaller and not consistently distinguishable from zero. Reporting increases are visible across most parties, again with the largest effects for AfD content.

We next examine whether treatment effects depend on respondents’ partisan alignment with the sender. As shown in Figure~\ref{fig:patterns-partisanship}, reductions in liking are strongest when emotionally manipulative content is posted by a respondent’s preferred party. This pattern suggests that the workshop reduced engagement even with politically congenial content, consistent with interrupting identity-consistent reinforcement of emotionally manipulative appeals. We do not detect comparably robust heterogeneity for sharing behavior.

\subsection*{Attention and exposure: skipping and watch time}

We then assess whether reduced liking reflects disengagement or more deliberate processing. As shown in Figure~\ref{fig:overall-patterns_other} in the SM, treated respondents were less likely to skip videos and spent a larger proportion of each video’s duration watching content. These patterns hold across political and entertainment videos. The combination of reduced liking of emotionally manipulative political posts and increased attention is consistent with more reflective engagement rather than avoidance.


\subsection*{Long-term effects}\label{longterm-effects}

To assess durability, we re-exposed respondents in a preregistered follow-up wave to a new set of political and entertainment videos (Figure~5). The follow-up was fielded from March~14 to April~4, 2025 (approximately 3--6 weeks after the first wave, depending on interview date).

For \textit{liking} and \textit{sharing}, treatment effects remained directionally consistent with the immediate post-treatment results, but attenuated in magnitude; estimates were smaller than in the first wave and no longer statistically distinguishable from zero for liking. For \textit{reporting} and \textit{bookmarking}, point estimates reversed sign relative to the short-run findings. 

Since  statistical power is lower due to panel attrition, we treat these longer-term results as exploratory. Attrition analyses show no systematic relationship between treatment assignment and dropout (SM, section~\ref{attrit}), although the smaller panel reduces precision. These findings therefore suggest that the behavioral consequences of digital literacy interventions may evolve over time, highlighting the importance of studying not only immediate, but also dynamic effects.




\subsection*{Downstream effect on AfD support}
Finally, we examine whether behavioral changes translated into downstream political preferences. This analysis was not preregistered and should therefore be interpreted as exploratory. We estimate treatment effects on within-person changes in AfD propensity to vote (PTV, 0--10) from baseline to each post-treatment wave (see section \ref{afd_support} in the SM).

Immediately after the intervention (wave~1), the estimated effect is negative, but not statistically distinguishable from zero: the workshop reduces the change in AfD PTV by −0.06 points on the 0--10 scale. Although imprecisely estimated, this corresponds to a non-trivial shift in relative terms, for example, roughly five to six out of 100 respondents would decrease their AfD PTV by one scale point compared to the control group, with most respondents unchanged. 

Three weeks later (wave~2), we detect a smaller, but statistically significant negative effect on AfD PTV change (−0.02 points on the 0--10 scale). Substantively, this corresponds to about two out of 100 respondents reducing their AfD PTV by one point relative to the control group.

Attrition analyses (section~\ref{attrit} in the SM) indicate that differential dropout does not account for this pattern. To interpret the persistence, we examine changes in party identification between waves (section~\ref{afd_change} in the SM). AfD affinity increased more strongly in the control group (+2.14 percentage points) than in the treatment group (+0.39 percentage points), consistent with contemporaneous increases in public support for the AfD during the field period (from 20\% to 23\%; \cite{PoliticoPolls}. Taken together, these results suggest that the intervention did not generate large immediate preference shifts, but may have modestly dampened responsiveness to contemporaneous increases in AfD support.


%\subsection*{Content Evaluation}

%\subsubsection*{Overall patterns}\label{overall-patterns}
%We begin by examining whether the workshop altered how individuals responded to political videos. As shown in Figure~1, participants in the intervention group engaged differently: they were roughly one percentage point less likely to \textit{like} emotionally manipulative political posts than those in the control group. For neutral political and entertainment content, treatment effects were small and not statistically significant. 

%The intervention also affected \textit{reporting} behavior. Participants in the treatment group were 1.3 percentage points more likely to report emotionally manipulative political posts than those in the control group, a sizeable relative increase given that only 2 percent of all videos were reported. Reporting of neutral or entertainment content remained unchanged. 

%In sum, the workshop enhanced participants’ ability to recognize and resist emotional manipulation in political videos. However, we find no significant effect on sharing behavior. Treated participants shared such posts slightly less often, but the difference was not statistically meaningful.

%Bookmarking increased by about 0.6 percentage points for political posts. Because bookmarking can reflect heterogeneous motives (e.g., saving content to revisit rather than endorsement), we interpret this outcome cautiously and focus on likes and reports as clearer evaluative and governance-related responses.


%\subsubsection*{Patterns by Party}\label{patterns-party}
%Beyond the overall effects, we examine whether treatment effects vary by the political party disseminating the videos. Figure~2 shows marked heterogeneity. The strongest effect appears for the AfD which is classified as a populist far-right party \cite{RooduijnEtAl2023_PopuListArticle,RooduijnEtAl2023_PopuListDataset}: participants in the treatment group were more than 2~percentage points less likely to like its manipulative posts compared to the control group. Similar though weaker trends are visible for other parties, with a statistically significant reduction only for \textit{Die Linke} which is classified as borderline populist and borderline far-left  \cite{RooduijnEtAl2023_PopuListArticle,RooduijnEtAl2023_PopuListDataset}. This suggests that the workshop curbed liking behavior primarily for manipulative content from both ideological extremes.

%\textit{Reporting} shows a different pattern. Treated participants were more likely to report manipulative posts from nearly all parties. The increase was largest for the AfD, but statistically significant for every party except the newly founded populist BSW. For some parties, reporting also rose for neutral political posts, indicating a broader shift in sensitivity to political content.

%By contrast, \textit{sharing} behavior did not differ by party. The increase in \textit{bookmarking}, however, appears concentrated on posts from the Christan democratic CDU/CSU and the liberal FDP, the only parties for which we find significant effects.

%\subsubsection*{Patterns by Strategy}\label{patterns-strategy}
%To assess whether the effects of the workshop vary depending on the manipulation strategy employed in political videos, we stratified the analysis of all emotionally manipulative political posts by strategy type, namely, fear, anger, and superiority. As shown in Figure~3, posts that reference superiority or anger are notably less likely to be liked and more likely to be reported by workshop participants compared to non-participants. Interestingly, the reduction in liking behavior among participants is observed only for posts employing the superiority or anger strategy; for posts invoking fear, we do not find statistically significant treatment effects with respect to liking behavior. In contrast, reporting behavior appears to be influenced by the workshop across all three strategy types, with the strongest effect observed for videos using the superiority strategy. At the same time, there do not seem to be strategy-specific effects with respect to sharing and bookmarking.







\section*{Conclusion}
Our study provides causal evidence that even brief online digital literacy workshops can reduce citizens’ susceptibility to emotionally manipulative political content in algorithmic video feeds. Participants exposed to the intervention were less likely to like and more likely to report such posts, especially those using anger or superiority appeals, disseminated by populist parties, or aligned with participants’ own preferences. At the same time, they skipped less and watched longer, suggesting more reflective engagement rather than avoidance. 
The combination of reduced positive reinforcement, increased reporting, and longer viewing times suggests that the intervention fostered more reflective engagement rather than simple disengagement.

These results demonstrate that scalable digital literacy interventions can work in precisely the environments where populist communication is most pervasive. The findings are especially relevant given that young people, who are the most active users of short-form video platforms, are also a key target group for populist parties. As political discourse increasingly migrates to algorithmic platforms, fostering emotional discernment and digital resilience may be critical to safeguarding democratic deliberation.

At the same time, our estimates capture immediate individual-level engagement under controlled exposure and do not incorporate algorithmic personalization or downstream diffusion in social networks. Future work should test how similar interventions perform under platform-level targeting and over longer time horizons.




%%%% NEW 
\section*{Materials and Methods}

\textbf{Preregistration, ethics, and sample}
The study was preregistered at OSF (\url{https://osf.io/js9ra/}) and approved by the ethics committee at Humboldt University of Berlin (HU-KSBF-EK\_2024\_0045). We recruited $N=4{,}437$ respondents from an online panel provider (Bilendi) using quotas for age, gender, and region. Eligibility required active use of short-form video platforms (e.g., TikTok, Instagram Reels, YouTube Shorts). Wave~1 was fielded from February~17 to February~23, 2025; a follow-up wave was fielded from March~14 to April~4, 2025.  We conducted the study in the context of the German election in 2025 to test the intervention in a high-stakes environment. 

\textbf{Design and intervention}
Our intervention adapts the approach developed by Athey et al. \cite{carletonathey2023} to target the role of emotional manipulation in political communication on social media. Participants were randomly assigned to either a control group or a treatment group that received the intervention. Participation in the intervention workshop was compulsory, with no option to opt out. Delivered as a brief online workshop in the context of the survey, the intervention aims to equip participants with the skills to recognize and resist emotional triggers, such as fear, anger, or superiority, that are frequently used in political videos to drive engagement and influence behavior. The workshop combines brief instructional modules with interactive exercises to help participants (1) identify emotionally charged content, (2) understand how emotional appeals can shape perceptions and sharing behavior, and (3) apply simple “stop-and-think” strategies when confronted with emotionally provocative posts. By promoting emotional awareness and critical reflection, the intervention seeks to reduce impulsive engagement with manipulative political content and foster more intentional online behavior.  Full intervention materials are provided in the Supplementary Materials (Section~\ref{appendix:intervention}).

\textbf{Stimuli and exposure environment}
After the intervention (treatment) or immediately after baseline items (control), respondents interacted for five minutes with a custom-built web platform designed to mimic core TikTok affordances (single-video display; standard engagement buttons; sequential exposure; see SM, section~\ref{video_engagement}). Each respondent was shown a randomized sequence of up to 400 short videos drawn from a fixed pool of 400 stimuli comprising 219 political videos and 181 non-political entertainment videos. Political videos were collected from official German party TikTok accounts posted during the 2025 federal election campaign period (see SM, section~\ref{videos}). Entertainment videos were collected from non-political “For You” feeds and screened to exclude political or offensive content (see SM, section~\ref{videos}). Videos were presented one at a time in random order; respondents could like, share, bookmark, report, re-watch, or skip to the next video. All interactions and watch-time were recorded unobtrusively. To prevent respondents from getting stuck on a few long videos, we restricted political videos to a maximum length of 40 seconds and entertainment videos to 20 seconds. 


\textbf{Videos}
Party videos were collected from a comprehensive list of 391 TikTok user handles manually identified by a research assistant (we scraped videos from 295 political users). We scraped the profile pages of each political TikTok party profile, as well as the videos and associated metadata, between September~1, 2024, and January~31, 2025, to capture the most relevant and recent political videos. Our video data includes all videos posted by these accounts up to that date, excluding videos labeled as "18 and older" by TikTok, as accessing them would have required logging into a real account. Entertainment content was collected by scraping the general TikTok “For You” page using a number of different German VPN locations to emulate the videos most likely to show up if scrolling the platform as a new user in Germany. These videos cover the period between October 2024 and the end of January 2025. Our dataset consisted of 2,793 political videos from 295 user accounts across the seven party categories and 743 general entertainment videos. 

\textbf{Coding of emotional manipulation}
Political videos were labeled by multiple independent coders for emotional manipulation strategies (fear, anger, superiority) using a pre-specified codebook. We retained videos with high agreement and excluded ambiguous cases; intercoder reliability and the full coding protocol are reported in the SM (section~\ref{videos} and section~\ref{appendix:codebook}). 

\textbf{Outcome measures}
Our preregistered behavioral outcomes capture engagement with political videos on the platform. For each respondent, we computed the proportion of emotionally manipulative political videos that were (i) liked, (ii) shared, (iii) reported, and (iv) bookmarked, as well as viewing outcomes including (v) the proportion of videos skipped and (vi) average watch time relative to video length. Parallel measures were computed for neutral political and entertainment content as placebo/benchmark outcomes (details in SM, Outcome Measurement section). In the follow-up wave, respondents were re-exposed to a new set of videos and the same behavioral outcomes were collected; follow-up analyses are reported as exploratory due to lower power and because they were not preregistered.

\textbf{Estimation}
We estimate average treatment effects by comparing treated and control respondents’ behavioral outcomes. For video-level analyses, we fit linear probability models with video fixed effects and cluster-robust standard errors at the respondent level to account for within-respondent correlation. For respondent-level proportion outcomes, we use analogous linear models with heteroskedasticity-robust standard errors. All specifications and robustness checks are reported in the SM.

%\textbf{Data and code availability.}
%[ADD: repository link / conditions for access / anonymization statement / OSF materials, code, and analysis scripts.]



%We preregistered and tested the hypothesis that participants in the workshop intervention group would engage less with emotionally manipulative political posts (by sharing, liking, or bookmarking) compared to participants in the control group. Additionally, we examined whether participants in the intervention group were more likely to report emotionally manipulative political posts than those in the control group. 

%Our study adapts a successful intervention developed by Athey et al. \cite{carletonathey2023} - originally designed to reduce misinformation  via SMS in Kenya - and applies it to political video content on TikTok. While Athey et al. \cite{carletonathey2023} demonstrated that emotion-based interventions outperformed reasoning-based ones in reducing impulsive sharing, their study relied on scripted text messages. By contrast, our study tests the effectiveness of a similar intervention in a much noisier real-world setting: the TikTok platform, using real political videos that parties released on TikTok in the 2025 German federal election campaign.















%%%%%% Figures and tables: These should be inserted at the end of the main text, as below (not in the middle of the text). Refer to them using e.g. Figure 1 (or Fig. 1) and Table 1.

\begin{figure}[ht]
\centering
\includegraphics[width=0.95\linewidth]{skeleton_files/figure-pdf/f_I.pdf}

\caption{Overall patterns}
\label{fig:overall-patterns}
\end{figure}

\begin{figure}
\centering
\includegraphics[width=0.90\linewidth]{skeleton_files/figure-pdf/f_IIa.pdf}
\caption{Patterns conditional on sender}
\label{fig:patterns-party}
\end{figure}


\begin{figure}[ht]
\centering
\includegraphics[width=0.95\linewidth]{skeleton_files/figure-pdf/f_IIc.pdf}
\caption{Patterns conditional on partisanship}
\label{fig:patterns-partisanship}
\end{figure}


\begin{figure}[ht]
\centering
\includegraphics[width=0.95\linewidth]{skeleton_files/figure-pdf/f_I_other.pdf}
\caption{Skipping and watch time}
\label{fig:overall-patterns_other}
\end{figure}


\begin{figure}[ht]
\centering
\includegraphics[width=\linewidth]{skeleton_files/Stata Graphs/f_I.pdf}
\caption{Long-term effects}
\label{fig:long}
\end{figure}


%%% REFERENCES
\clearpage

\bibliography{scibib}
\bibliographystyle{Science}









\clearpage
\begin{center}
\appendix{\textbf{\Large{Supplementary Materials}}}
\end{center}
\appendix
\counterwithin{figure}{section}
\counterwithin{table}{section}
\tableofcontents
\clearpage





\section{Descriptives}
\begin{table}[h]
    \caption{Descriptives on videos and video behavior \label{tab:descriptives}}
    \centering
    {\footnotesize \def\arraystretch{0.95}
    \begin{tabular}{p{0.5cm} p{6.5cm} r r p{0.5cm}}
    \hline
        \input{skeleton_files/tables/descriptives_videos} \\
    \hline
    \end{tabular}
    }

    \vspace{0.3cm}
    \parbox{0.9\linewidth}{\footnotesize \textit{Notes: The short-term wave took place between February 17, 2025, and February 23, 2025, while the long-term wave took place between March 14, 2025, and April 4, 2025.}}
\end{table}


%%%%%%%%%%%%%%%%%%%
\clearpage
\section{Scraping of TikTok videos}\label{tiktok:users_list}
Below, we include a table containing all the political user handles for which we collected videos and metadata. We crawled the TikTok pages of the 391 political user profiles originally identified and retrieved profile and video metadata for 367. Using the collected metadata, we narrowed down the videos to our chosen time window (September 2024 - January 2025) and the duration limit of 40 seconds. Following these steps, we were left with videos collected from 295 political users. Of these, 103 were included in the video feed sample shown to respondents (as indicated by the "In experiment" column). 
\input{skeleton_files/tables/political_userhandles_scraped}





\clearpage
\section{Coding of TikTok Videos}
\label{videos}

\begin{table}[ht]
    \centering
\begin{tabular}{|r|r|r|}
\hline
\textbf{Anger} & \textbf{Fear} & \textbf{Superiority} \\
\hline
0.469 & 0.457 & 0.376 \\
\hline
\end{tabular}
    \caption{Average Cohen's Kappa between coder pairings}
    \label{tab:cohens_kappa}
\end{table}

\begin{table}[ht]
\centering
\begin{tabular}{|l|l|l|l|l|l|}
\hline
                              & \textbf{anger}             & \textbf{fear}              & \textbf{superiority}       & \textbf{neutral}            & \textbf{all} \\ \hline
AfD                           & 9                          & 9                          & 9                          & 9                           & 36           \\ \hline
BSW                           & 1                          & 2                          & 4                          & 9                           & 16           \\ \hline
CDU/CSU                       & 9                          & 9                          & 9                          & 9                           & 36           \\ \hline
FDP                           & 9                          & 5                          & 9                          & 9                           & 32           \\ \hline
Grüne                         & 9                          & 7                          & 9                          & 9                           & 34           \\ \hline
Linke                         & 9                          & 4                          & 7                          & 9                           & 29           \\ \hline
SPD                           & 9                          & 9                          & 9                          & 9                           & 36           \\ \hline
Entertainment                 & 0                          & 0                          & 0                          & 181                         & 181          \\ \hline
Total                         & 55                         & 45                         & 56                         & 244                         & 400          \\ \hline
\end{tabular}
\caption{Strategy distribution across videos}
\label{tiktok:strategy}
\end{table}

Videos for manual coding were sampled randomly, but proportionally across each party. To ensure that our entertainment videos were truly non-political, coders manually combed through random samples of the content for political relevance, as well as for potentially offensive or inappropriate material. For political videos, coders used a codebook and engaged in iterative discussions with the team and each other. Videos were coded according to predefined categories of political emotional manipulation: fear, anger, superiority, and neutral. A video could be assigned multiple categories if the coder deemed it appropriate. Each coder received a sample of videos spanning all seven parties. We conducted multiple rounds of coding, ensuring that each video was coded by at least two different coders. Additionally, a third round was conducted in which a proportion of the videos were coded by an additional coder. This resulted in approximately 1,327 coded TikTok videos across parties.  

To ensure reliable labeling across videos, we followed a structured validation process: (i) For videos coded by three coders, we retained the label agreed upon by at least two coders, (ii) for videos coded by only two coders, we discarded videos where there was disagreement and kept only those with matching labels, (iii) to maximize clarity in the experiment, we further excluded videos coded with multiple emotional categories, resulting in a narrower but more clear-cut sample. The intercoder reliability was moderate (see table ~\ref{tab:cohens_kappa}). This was partially due to differing perceptions among coders regarding the intensity of fear, anger, or superiority expressed in the videos, leading to variation in when coders considered a particular category sufficiently present.

The final subset of manually coded data was sampled into a set of 400 videos for the experiment, consisting of a random sample of 181 entertainment videos and 219 political videos. Among our political videos, we randomly sampled nine videos per strategy for each party. Some parties with fewer videos available were sampled fully (BSW). Table ~\ref{tiktok:strategy} provides an overview of the final 400 videos, broken down by strategy and party. 
The final videos shown to respondents included the original video, user handle, profile image, description, and hashtags, maintaining the authentic look and feel of TikTok content.






\clearpage




\section{Codebook for coding of videos}\label{appendix:codebook}

\textbf{Practical instructions} \\
You have received an Excel file linked to your randomly generated RA number with a random sample of TikTok video URLs across seven party categories (AfD, BSW, CDU/CSU, FDP, Grüne, Linke, and SPD).

The file further includes three columns: fear, anger, and superiority. If a category is present in a video, the cell of the column should be denoted as 1, else as 0. One TikTok video can be marked with more than one of the three categories, as well as zero across all three.

You should watch each video in its entirety (no video is longer than 40 seconds) and code in accordance with all aspects of the video content, including speech, imagery, sounds, etc.


\textbf{Fear} \\
Content that elicits fear, anxiety, or concern by focusing on risks, dangers, and threats to safety and/or values - messages emphasizing uncertainty and potential personal or societal harm (including financial insecurity).

\textit{Indicators:}
\begin{itemize}
    \item Warnings of imminent harm or disaster.
    \item Use of alarming statistics or dramatic imagery.
    \item Language that emphasizes urgency or helplessness (e.g., ``crisis,'' ``catastrophe,'' ``you're at risk'').
\end{itemize}

\textit{Examples:}
\begin{itemize}
    \item ``If we don't act now, everything will be lost.''
    \item ``They are coming for your freedom.''
\end{itemize}



\textbf{Anger}\\
Content aimed at provoking anger or outrage by highlighting (perceived) injustices, immorality, or wrongs, including calls to action and portrayals of protest and resistance.

\textit{Indicators:}
\begin{itemize}
    \item Accusatory language targeting individuals, groups, or institutions.
    \item Emphasis on betrayal, unfairness, or corruption.
    \item Use of inflammatory words (e.g., ``traitor,'' ``liar,'' ``scam'').
\end{itemize}

\textit{Examples:}
\begin{itemize}
    \item ``This is a betrayal of our people.''
    \item ``They are stealing from us, and no one is stopping them!''
    \item ``They are destroying our country!''
\end{itemize}



\textbf{Superiority}\\
Content that seeks to evoke feelings of moral or intellectual superiority over others.

\textit{Indicators:}
\begin{itemize}
    \item Depictions of ``us versus them'' dynamics, emphasizing the moral or intellectual flaws of others.
    \item Appeals to elitism or exclusivity (e.g., ``Only the smart ones will understand'') as well as populism (e.g., ``those big city elites will never understand our struggles'').
    \item Use of belittling or dismissive language toward opposing viewpoints.
\end{itemize}

\textit{Examples:}
\begin{itemize}
    \item ``We know the truth; they're just blind followers.''
    \item ``This is what sets us apart from the rest.''
\end{itemize}









\clearpage
\section{Intervention}\label{appendix:intervention}


\subsection*{General introduction}
\includegraphics[width=0.95\linewidth]{skeleton_files/screenshots/Workshop_01.png}\par\vspace{2ex}
\includegraphics[width=0.95\linewidth]{skeleton_files/screenshots/Workshop_02.png}\par\vspace{2ex}
\includegraphics[width=0.95\linewidth]{skeleton_files/screenshots/Workshop_03.png}\par\vspace{2ex}
\includegraphics[width=0.95\linewidth]{skeleton_files/screenshots/Workshop_04.png}\par\vspace{2ex}
\includegraphics[width=0.95\linewidth]{skeleton_files/screenshots/Workshop_05.png}\par\vspace{2ex}

\subsection*{Fear}
\includegraphics[width=0.95\linewidth]{skeleton_files/screenshots/Workshop_06.png}\par\vspace{2ex}
\includegraphics[width=0.95\linewidth]{skeleton_files/screenshots/Workshop_07.png}\par\vspace{2ex}
\includegraphics[width=0.95\linewidth]{skeleton_files/screenshots/Workshop_08.png}\par\vspace{2ex}
\includegraphics[width=0.95\linewidth]{skeleton_files/screenshots/Workshop_09.png}\par\vspace{2ex}
\includegraphics[width=0.95\linewidth]{skeleton_files/screenshots/Workshop_10.png}\par\vspace{2ex}

\subsection*{Anger}
\includegraphics[width=0.95\linewidth]{skeleton_files/screenshots/Workshop_11.png}\par\vspace{2ex}
\includegraphics[width=0.95\linewidth]{skeleton_files/screenshots/Workshop_12.png}\par\vspace{2ex}
\includegraphics[width=0.95\linewidth]{skeleton_files/screenshots/Workshop_13.png}\par\vspace{2ex}
\includegraphics[width=0.95\linewidth]{skeleton_files/screenshots/Workshop_14.png}\par\vspace{2ex}
\includegraphics[width=0.95\linewidth]{skeleton_files/screenshots/Workshop_15.png}\par\vspace{2ex}

\subsection*{Superiority}
\includegraphics[width=0.95\linewidth]{skeleton_files/screenshots/Workshop_16.png}\par\vspace{2ex}
\includegraphics[width=0.95\linewidth]{skeleton_files/screenshots/Workshop_17.png}\par\vspace{2ex}
\includegraphics[width=0.95\linewidth]{skeleton_files/screenshots/Workshop_18.png}\par\vspace{2ex}
\includegraphics[width=0.95\linewidth]{skeleton_files/screenshots/Workshop_19.png}\par\vspace{2ex}
\includegraphics[width=0.95\linewidth]{skeleton_files/screenshots/Workshop_20.png}\par\vspace{2ex}
\includegraphics[width=0.95\linewidth]{skeleton_files/screenshots/Workshop_21.png}\par\vspace{2ex}
\includegraphics[width=0.95\linewidth]{skeleton_files/screenshots/Workshop_22.png}\par\vspace{2ex}

\vspace{0.3cm}

\textbf{Literature:}
\begin{enumerate}
    \item \label{ryan2012click} Ryan, T. J. (2012). \textit{What Makes Us Click? Demonstrating Incentives for Angry Discourse with Digital-Age Field Experiments}. The Journal of Politics, 74(4), 1138–1152.
    \item \label{walter2019meta} Walter, Nathan, et al. (2019). \textit{Meta-Analysis of Anger and Persuasion: An Empirical Integration of Four Models}. Journal of Communication, 69(1), 73–93. \\https://doi.org/10.1093/joc/jqy054
\end{enumerate}






%--------------------------------------------
\clearpage
\section{Outcome Measurement}
After the experimental block, participants were told that the following part of the study concerned the use of social media platforms. Participants were informed that they would now spend five minutes on a “new online video platform” that functions very similarly to TikTok. They were asked to imagine that they had just signed up to this platform for the first time and were given five minutes to browse the feed and react to the videos.

Before the videos appeared, participants received a brief explanation of the interface. They were told that the goal of the study was to understand how people use this new platform and respond to videos “as they would on a real social media platform,” and were encouraged to behave naturally. The interface offered standard TikTok-style engagement options: like (heart icon), share, bookmark, report (for problematic videos), and skip to the next video. After the five-minute viewing period, participants were thanked, asked whether all videos had displayed correctly, and could describe any technical problems or provide additional comments about the platform in open-ended text fields.

\clearpage
\subsection*{Screenshots}
"In the following, we are interested in the use of social media platforms on the internet."
\includegraphics[width=0.95\linewidth]{skeleton_files/screenshots/Videos_Intro_1.png}\par\vspace{2ex}
"As in our survey a few weeks ago, we are again interested today in the use of video platforms. In a moment you will again spend five minutes on a video platform. This video platform is a new platform that works very similarly to TikTok."

\includegraphics[width=0.95\linewidth]{skeleton_files/screenshots/Videos_Intro_2.png}\par\vspace{2ex}
"Now imagine that you have just signed up for the first time on the new online video platform. You now have five minutes to look around this new platform and respond to the videos that are shown."

\includegraphics[width=0.95\linewidth]{skeleton_files/screenshots/Videos_Intro_3.png}\par\vspace{2ex}

\clearpage
"We would like to find out how people use the new online video platform and react to the videos – just as they would on a real social media platform. Please behave completely naturally and use the different buttons to express your preferences and interests as you normally would."

Buttons:
“Like: ‘Like’ the video”

“Share: Share the video with others”

“Bookmark: Save the video”

“Report video: Report a problematic video”

“Skip: Jump to the next video”

\includegraphics[width=0.95\linewidth]{skeleton_files/screenshots/Videos_Intro_4.png}\par\vspace{2ex}



\label{video_engagement}
\begin{figure}[ht]
\centering
\includegraphics[width=0.95\linewidth]{skeleton_files/screenshots/Videos_13.png}\caption{Example video}
\label{fig:video}
\end{figure}
%--------------------------------------------
\clearpage
\section{Ecological Validity of the Intervention}\label{appendix:ecological_validity}
The validity of our results depends on whether respondents engaged with our platform as they would with other social media platforms. This requires that they perceived the platform as real rather than part of a study. To assess the ecological validity of our setup and intervention, we asked respondents an open-ended question about the platform after they saw the videos and engaged with the platform: \textit{“Is there anything else you’d like to tell us about the video platform?”}

To analyze the data, we first cleaned the answers by removing cases where respondents did not provide any substantive comment. We then used a large language model (OpenAI \texttt{gpt-4o}, temperature = 0) to classify German free-text comments into four mutually exclusive categories indicating how respondents treated the platform and its content: 
\begin{itemize}
    \item \textbf{setup skepticism (0)} comments explicitly frame the platform as part of a study or as fake
    \item  \textbf{platform-focused (1)} comments treat the platform as real and discuss UX/features/algorithm or compare it to existing platforms such as TikTok/Instagram/YouTube—e.g., \emph{“dislike button is missing,” “scroll instead of click,” “like TikTok”} 
    \item \textbf{content-focused (2)} comments that treat the content as real—e.g., \emph{“too much politics,” “boring,” “waste of time,” “many fake news from parties”})
    \item \textbf{other/unclear (3)} comments that fall in the category are very short or vague remarks without a clear platform/content judgment—e.g., \emph{“ok,” “thanks,” “bad,” “good survey”}). 
\end{itemize}
Our prompt constrained the model to assign \textbf{1} when platform/UX terms or platform comparisons appeared (including platform name alone), and \textbf{2} when comments expressed positive/negative valence about the content or videos without questioning authenticity. When both platform and content were mentioned, the model was instructed to choose the primary focus (1 vs.\ 2). 


Figure \ref{fig:classification_comments} presents the results. The vast majority of respondents engaged either with the design of the platform ($20\%$) or the nature of the content on the platform ($46\%$). A large number of respondents ($34\%$) also had general comments. Only a small minority ($1\%$) of comments were classified as setup skeptics. After manually checking comments that were classified as setup skeptic, we saw that just a handful of those actually perceived the platform as "fake" and part of an experiment. 

\begin{figure}[ht]
\centering
\includegraphics[width=0.95\linewidth]{skeleton_files/figure-pdf/fig_classification_comments.pdf}
\caption{Engagement with the platform}
\label{fig:classification_comments}
\end{figure}




%%%%%%%%%%%%%%%%%%%%%%%
\clearpage
\section{Further Analyses on Content Evaluation}\label{appendix:further_analyses}


\begin{figure}[ht]
\centering
\includegraphics[width=0.95\linewidth]{skeleton_files/figure-pdf/f_IIb.pdf}
\caption{Patterns conditional on strategy}
\label{fig:patterns-strategy}
\end{figure}


\begin{figure}[ht]
\centering
\includegraphics[width=0.95\linewidth]{skeleton_files/figure-pdf/f_IIc3.pdf}
\caption{Patterns conditional on partisanship and sender}
\label{fig:patterns-partisanship-party}
\end{figure}


\begin{figure}[ht]
\centering
\includegraphics[width=0.95\linewidth]{skeleton_files/figure-pdf/f_IIa_other.pdf}
\caption{Patterns conditional on sender}
\label{fig:patterns-party-other}
\end{figure}

\begin{figure}[ht]
\centering
\includegraphics[width=0.95\linewidth]{skeleton_files/figure-pdf/f_IIb_other.pdf}
\caption{Patterns conditional on strategy}
\label{fig:patterns-strategy-other}
\end{figure}

\begin{figure}[ht]
\centering
\includegraphics[width=0.95\linewidth]{skeleton_files/figure-pdf/f_IIc_other.pdf}
\caption{Patterns conditional on partisanship}
\label{fig:patterns-partisanship-other}
\end{figure}

\begin{figure}[ht]
\centering
\includegraphics[width=0.95\linewidth]{skeleton_files/figure-pdf/f_IIc3_other.pdf}
\caption{Patterns conditional on partisanship and sender}
\label{fig:patterns-partisanship-party-other}
\end{figure}





%%%%%%%%%%%%%%%%%%%%%%%
\clearpage
\section{Attrition}
\label{attrit}
Since wave 2 was fielded a few weeks after wave 1, we assess potential differential attrition across the treatment and control group. We define an indicator
$\textit{attrit}_i=1$ if respondent $i$ is missing in wave 2 and $0$ otherwise, and estimate separate linear probability models
\begin{equation}
\textit{attrit}_i \;=\; \alpha \;+\; \beta_1\,\textit{treated}_{i,w9} \;+\; \beta_2\,X_i \;+\; \beta_3\big(\textit{treated}_{i,w9}\times X_i\big) \;+\; \varepsilon_i,
\end{equation}
where $X$ cycles through pre-treatment covariates (age, gender, region/state dummies, urbanization, social-media use, political
interest, ideology, party ID and pre-treatment PTVs, socio-economic variables, etc.). We report the interaction coefficient
$\beta_3$ (difference in attrition between treatment and control across values of $X$) with 95\% confidence intervals. To maximize power, each regression uses the largest available sample complete on
$\{\textit{attrit},\,\textit{treated}_{w9},\,X\}$.
In each case we standardize $X$ to have unit standard deviation. 

\begin{figure}[ht]
\centering
\includegraphics[width=0.95\linewidth]{skeleton_files/figure-pdf/fig_attrit_interaction.pdf}
\caption{Attrition (0/1) on Treatment-Covariate Interaction}
\label{fig:fig_attrit_interaction}
\end{figure}





%%%%%%%%%%%%%%%%%%%%%%%
\clearpage
\section{Downstream effects on AfD support}
\label{afd_support}


\begin{figure}[ht]
\centering
\includegraphics[width=0.95\linewidth]{skeleton_files/figure-pdf/afd_short.png}
\caption{Effect on PTV (first wave)}
\end{figure}


\begin{figure}[ht]
\centering
\includegraphics[width=0.95\linewidth]{skeleton_files/figure-pdf/afd_long.png}
\caption{Effect on PTV (second wave)}
\end{figure}




%%%%%%%%%%%%%%%%%%%%%%%
\clearpage
\section{Changes in Party Identification from Wave 1 to Wave 2}
\label{afd_change}
Below, we plot the change in party identification from wave 1 to wave 2 for the full sample and separated by treatment and control group. 

\begin{figure}[ht]
\centering
\includegraphics[width=0.95\linewidth]{skeleton_files/figure-pdf/fig_pid_change.pdf}
\caption{Change in Party Identification, Wave 1 to Wave 2, Full Sample}
\end{figure}


\begin{figure}[ht]
\centering
\includegraphics[width=0.95\linewidth]{skeleton_files/figure-pdf/fig_pid_change_treat.pdf}
\caption{Change in Party Identification, Wave 1 to Wave 2, Treatment Group}
\end{figure}


\begin{figure}[ht]
\centering
\includegraphics[width=0.95\linewidth]{skeleton_files/figure-pdf/fig_pid_change_control.pdf}
\caption{Change in Party Identification, Wave 1 to Wave 2, Control Group}
\end{figure}


\clearpage 
\section{Correlation between short-term and long-term outcomes }

\begin{figure}[ht]
\centering
\includegraphics[width=0.95\linewidth]{skeleton_files/correlation_ellipses.pdf}
\caption{Correlation between short-term and long-term outcomes}
\label{fig:correlation}
\end{figure}



\clearpage 
\section{Reconciliation}

Table~\ref{tab:reconciliation} summarizes key deviations between our preregistered pre-analysis plan (PAP) and the analyses reported in the manuscript.

\begin{table}[ht]
\centering
\caption{Reconciliation between the pre-analysis plan (PAP) and the manuscript}
\label{tab:reconciliation}
\small
\begin{tabular}{p{3.2cm} p{5.8cm} p{5.8cm}}
\toprule
\textbf{Item} & \textbf{PAP} & \textbf{Manuscript} \\
\midrule
Video-feed randomization & Planned a constrained feed (max.\ five political videos in a row) with a balancing procedure; PAP documents a malfunction during early fieldwork and a corrective switch to unconstrained random draws. & Results use the corrected implementation: sequential exposure with an unconstrained random draw from a fixed stimulus pool. \\
Primary behavioral outcomes & Focus on engagement with emotional political posts (share/like/bookmark) and reporting; long-term persistence posed as H2; skipping/watch time listed as exploratory.  & Emphasizes \textit{liking} and \textit{reporting} as primary behavioral outcomes in the main text; attention outcomes (skipping/watch time) and heterogeneity are presented centrally as evidence on processing. \\
Follow-up (wave 2) analyses & Follow-up wave planned to assess longer-term effects (H2)  & Wave-2 results are reported as exploratory due to lower power and attrition  \\
Downstream political preferences & Not included as a preregistered outcome in the PAP.  & Adds exploratory analysis of downstream AfD support (PTV change) and related mechanisms (party identification changes). \\
\bottomrule
\end{tabular}
\end{table}

\end{document}